\subsection{Primitive ideals}

Let A be an algebra over K and E a vector space over K. One calls a representation of A in E a morphism from A to the algebra \(\mathcal{L}(\mathrm{E})\) of endomorphisms of E. An injective representation is said to be faithful. Let \(\pi_1\) and \(\pi_2\) be representations of A in spaces \(\mathrm{E}_1\) , \(\mathrm{E}_2\) . A morphism of \(\pi_1\) into \(\pi_2\) is a K-linear map \(u: \mathrm{E}_1 \to \mathrm{E}_2\) such that \(u(\pi_1(a)x) = \pi_2(a)u(x)\) for all \(a \in \mathrm{A}\) and \(x \in \mathrm{E}_1\) . Representations are said to be equivalent if there exists a morphism of \(\pi_1\) into \(\pi_2\) which is an isomorphism of vector spaces. Its inverse is then a morphism of \(\pi_2\) into \(\pi_1\) . A representation \(\pi\) of A in E is said to be irreducible if \(\mathrm{E} \neq \{0\}\) and if the only vector subspaces of E stable under \(\pi(\mathrm{A})\) are \(\{0\}\) and E.

Example. --- The zero map from A to \(\mathcal{L}(\mathrm{E})\) is a representation, called the trivial representation, of A. It is irreducible if and only if E has dimension 1.

Lemma 4. --- Let π be an irreducible, non-trivial representation of A in E. For every non-zero element ξ of E, one has π(A)ξ = E.

The subspace \(\pi(A)\xi\) of \(E\) is stable under \(\pi(A)\) . Suppose it is zero. The non-zero subspace \(K\xi\) of \(E\) would then be stable under \(\pi(A)\) , and hence equal to \(E\) ; but this would imply that \(\pi\) is the zero representation. We therefore have \(\pi(A)\xi = E\) .

Let \(\pi\) be an irreducible, non-trivial representation of A in E. By this lemma, the annihilator R of \(\xi\) in A is a regular left ideal (A, VIII, p.~425, n° 1) of A, and the representation \(\pi\) is equivalent to the representation defined by the A-pseudomodule A/R. Since \(\pi\) is irreducible, the ideal R is a maximal regular left ideal.

DEFINITION 7. --- Let A be an algebra over K. One calls a primitive ideal of A the kernel of a non-trivial irreducible representation of A.

If A is commutative, the primitive ideals of A are the maximal regular ideals of A. Indeed, the non-trivial irreducible representations of A are, up to equivalence, the representations \(\pi_{R}\) defined by the A-pseudomodules A/R, where R is a maximal regular ideal of A. The kernel of \(\pi_{R}\) contains R. It is even equal to it since, by A, VIII, p.~426, prop. 2, the commutativity of A implies that A/R is a field. Hence \(\text{Ker}(\pi_{\text{R}})\) is maximal regular.

Lemma 5. --- Let π be an irreducible representation of A in a vector space E over K.

\begin{enumerate}
\def\labelenumi{\alph{enumi})}
\item
  Let I be a two-sided ideal of A. If \(\pi(\mathrm{I}) \neq \{0\}\) , then \(\pi|\mathrm{I}\) is irreducible;
\item
  Let \(I_{1}\) and \(I_{2}\) be two-sided ideals of A such that \(\pi(\mathrm{I}_{1}) \neq 0\) and \(\pi(\mathrm{I}_{2}) \neq 0\) . Then \(\pi(\mathrm{I}_{1}\mathrm{I}_{2}) \neq 0\) .
\end{enumerate}

The set of elements of E annihilated by \(\pi(\mathrm{I})\) is stable under \(\pi(\mathrm{A})\) and distinct from E, hence equal to 0. Therefore, if \(\xi\) is a non-zero element of E, one has \(\pi(\mathrm{I})\xi \neq 0\) ; since \(\pi(\mathrm{I})\xi\) is stable under \(\pi(\mathrm{A})\) , one has \(\pi(\mathrm{I})\xi = \mathrm{E}\) , which proves \(a\) ). On the other hand, what precedes proves that \(\pi(\mathrm{I}_2)\mathrm{E} = \mathrm{E}\) , \(\pi(\mathrm{I}_1)\pi(\mathrm{I}_2)\mathrm{E} = \mathrm{E}\) , hence \(\pi(\mathrm{I}_1\mathrm{I}_2) \neq 0\) , whence \(b\) ).

Lemma 6. --- Let I \(_{1}\) and I \(_{2}\) be two-sided ideals of A, and I a primitive ideal of A. If I contains I \(_{1}\) I \(_{2}\) (in particular, if I contains I \(_{1}\) ∩ I \(_{2}\) ), then I contains I \(_{1}\) or I \(_{2}\) .

Let \(\pi\) be an irreducible representation with kernel I. If I \(\not\supseteq\) I \(_1\) and I \(\not\supseteq\) I \(_2\) , Lemma 5, \(b\) ) proves that \(\pi(\text{I}_1\text{I}_2) \neq 0\) , whence I \(\not\supseteq\) I \(_1\) I \(_2\) .

Lemma 7. --- Suppose that A has a unit element. Let I be a maximal two-sided ideal of A. Then I is a primitive ideal.

There exists a maximal left ideal R of A containing I (A, I, p.~99, th. 1). Let \(\pi\) be the canonical representation of A in A/R, which is irreducible and nonzero. Since IA \(\subset\) R, the kernel I' of \(\pi\) contains I, hence I' = I and I is primitive.

Let J(A) be the set of primitive ideals of A. For any subset M of A, we denote by V(M) the set of primitive ideals of A containing M; if I is the two-sided ideal of A generated by M, then V(M) = V(I). If M is reduced to a single element x, we write V(x) instead of V(\{x\}).

The mapping M \(\mapsto\) V(M) is decreasing with respect to inclusion relations. We have:

\[
\begin{array}{r l} (5) & \mathrm{V} (\varnothing) = \mathrm{J} (\mathrm{A}), \qquad \mathrm{V} (\mathrm{A}) = \varnothing \end{array}
\]

\[
\begin{array}{r l} (6) & \mathrm{V} \left(\bigcup_ {i \in \mathrm{I}} \mathrm{M} _ {i}\right) = \mathrm{V} \left(\sum_ {i \in \mathrm{I}} \mathrm{M} _ {i}\right) = \bigcap_ {i \in \mathrm{I}} \mathrm{V} (\mathrm{M} _ {i}) \end{array}
\]

for any family \((\mathrm{M}_{i})_{i\in\mathrm{I}}\) of subsets of A. On the other hand, by Lemma 6,

\[
\begin{array}{r l} (7) & \mathrm{V} (\mathrm{I} _ {1} \cap \mathrm{I} _ {2}) = \mathrm{V} (\mathrm{I} _ {1} \mathrm{I} _ {2}) = \mathrm{V} (\mathrm{I} _ {1}) \cup \mathrm{V} (\mathrm{I} _ {2}) \end{array}
\]

for all two-sided ideals \(I_{1}\) , \(I_{2}\) of A. Formulas (5) through (7) show that the subsets V(M) of J(A) are the closed subsets of a topology called the Jacobson topology on J(A).

Let T be a subset of J(A), and let \(\Upsilon(T)\) be the intersection of the elements of T, so that \(\Upsilon(T)\) is a two-sided ideal of A. Then the closure of T in J(A) is the smallest closed subset of J(A) containing T, that is, V( \(\Upsilon(T)\) ). In particular, T is closed if and only if \(T = V(\Upsilon(T))\) .

PROPOSITION 2. --- Let I \(_{1}\) and I \(_{2}\) be distinct points of J(A). Then one of these two points is not adherent to the other.

Indeed, we have, for example, \(I_{1} \not\subset I_{2}\) . The set \(V(I_{1})\) of \(I \in J(A)\) such that \(I_{1} \subset I\) is closed in \(J(A)\) and it contains \(I_{1}\) but not \(I_{2}\) .

PROPOSITION 3. --- Let I ∈ J(A). For \{I\} to be closed in J(A), it is necessary and sufficient that I be a primitive maximal ideal.

Indeed, the closure of \{I\} consists of the primitive ideals of A containing I.

The relation

« \(\pi_{1},\pi_{2}\) are representations of A, which are isomorphic.

is an equivalence relation with respect to \(\pi_1\) and \(\pi_2\) . For any representation \(\pi\) of A, we denote by \(\operatorname{cl}(\pi)\) the equivalence class of \(\pi\) , which is therefore a representation of A isomorphic to \(\pi\) , such that two representations \(\pi_1\) and \(\pi_2\) are isomorphic if and only if \(\operatorname{cl}(\pi_1) = \operatorname{cl}(\pi_2)\) . We say that \(\operatorname{cl}(\pi)\) is the class of \(\pi\) .

Let c be the cardinal of A. Let \(\pi\) be a nonzero irreducible representation of A in a K-vector space E. Let \(\xi\) be a nonzero element of E. Since \(\pi(A)\xi = E\) (lemma 4), the dimension of E is \(\leqslant\mathfrak{c}\) (A, II, p.~97, corollary). The relation

“λ is a class of irreducible representations of A

in a K-vector space of dimension \(\leqslant\mathfrak{c}\gg\)

is collectivizing in \(\lambda\) (E, II, p.~3). Indeed, every vector space of dimension \(\leqslant \mathfrak{c}\) is isomorphic to a space \(\mathrm{K}^{\mathrm{B}}\) where B is a subset of A (A, II, p.~25, def. 10), and the assertion then follows from E, II, p.~47.

We denote by \(\widehat{\mathrm{A}}\) the set of classes of irreducible, nontrivial representations of A. According to what precedes, for every nontrivial irreducible representation \(\pi\) of A, there exists a unique representation \(\widehat{\pi} \in \widehat{\mathrm{A}}\) which is isomorphic to \(\pi\) .

The map from \(\hat{\mathrm{A}}\) into \(\mathrm{J(A)}\) which associates to \(\pi\) its kernel is surjective. If A is commutative, it follows from the fact that primitive ideals are regular maximal ideals that this map is a bijection.

We endow \(\hat{A}\) with the inverse-image topology under the map \(\hat{A} \rightarrow J(A)\) .

If A possesses a unit element, the spaces J(A) and \(\hat{\mathrm{A}}\) are quasi-compact.

It suffices to give the proof for J(A). Let \((\mathrm{T}_{j})\) be a family of closed subsets of J(A) whose intersection is empty. If the sum \(\sum_{j}\Upsilon(\mathrm{T}_{j})\) was different from A, then this sum would be contained in a maximal two-sided ideal I. The ideal I would be primitive (Lemma 7); since each \(T_{j}\) is closed, hence equal to \(\mathrm{V}(\Upsilon(\mathrm{T}_{j}))\) , we would have \(I\in T_{j}\) for all j, which contradicts the hypothesis. Thus we have \(\sum_{j}\Upsilon(\mathrm{T}_{j})=\mathrm{A}\) , and therefore we can write \(1=x_{1}+\cdots+x_{n}\) with \(n\geqslant1\) and \(x_{i}\in\Upsilon(\mathrm{T}_{j_{i}})\) for all i. This implies that \(\Upsilon(\mathrm{T}_{j_{1}})+\cdots+\Upsilon(\mathrm{T}_{j_{n}})=\mathrm{A}\) , whence \(T_{j_{1}}\cap\cdots\cap T_{j_{n}}=\varnothing\) .

Suppose the algebra A is commutative and unital. The Jacobson topology on J(A) is the topology induced on J(A) by the Zariski topology of the prime spectrum of A (AC, II, def. 4, p.~125).

Suppose that A is commutative and that K is a topological field. The canonical isomorphism from K onto \(\mathcal{L}(\mathrm{K})\) allows one to identify an element of \(\mathsf{X}(\mathsf{A})\) with a representation of A in the vector space K, which defines an injective map from \(\mathsf{X}(\mathsf{A})\) into \(\widehat{\mathsf{A}}\) . We can therefore identify \(\mathsf{X}(\mathsf{A})\) with a subset of \(\widehat{\mathsf{A}}\) .

Proposition 5. — The topology induced on \(\mathsf{X}(\mathsf{A})\) by that of \(\widehat{\mathsf{A}}\) is coarser than the topology of \(\mathsf{X}(\mathsf{A})\) .

Indeed, let T be a closed subset of \(\widehat{\mathbf{A}}\) . Then T is the set of \(\pi \in \widehat{\mathbf{A}}\) whose kernel contains a subset M of A. Therefore \(\mathrm{T} \cap \mathrm{X}(\mathrm{A})\) is the set of \(\chi \in \mathrm{X}(\mathrm{A})\) which vanish on M, that is, a closed subset of \(\mathrm{X}(\mathrm{A})\) . Whence the proposition.

In general, the topology of \(\mathsf{X}(\mathsf{A})\) does not coincide with the topology induced by the topology of \(\widehat{\mathsf{A}}\) (cf.~I, p.~193, exercise 6, c)).

